\uuid{TGd1}
\exo7id{2466}
\titre{exo7 2466}
\auteur{matexo1}
\organisation{exo7}
\datecreate{2002-02-01}
\isIndication{false}
\isCorrection{false}
\chapitre{Déterminant, système linéaire}
\sousChapitre{Système linéaire, rang}
\cpge{mpsi-mp2i.determinants.b, pcsi.determinants.a, ptsi.determinants.c, mpsi-mp2i.determinants.d, pcsi.calcul-matriciel-systemes-lineaires.d, ptsi.calcul-matriciel-systemes-lineaires.d, mpsi-mp2i.matrices-applications-lineaires.a, pcsi.matrices-applications-lineaires.a, ptsi.matrices-applications-lineaires.a, bcpst2.applications-lineaires-matrices.a}
\module{Algèbre}
\niveau{L2}
\difficulte{}

\contenu{
\texte{
Soit la matrice
$$A=\left( \begin{array}{ccc}
m-2 & 2 & -1 \\
2 & m & 2 \\
2m & 2m+2 & m+1
\end{array}\right).$$
}
\begin{enumerate}
    \item \question{Calculer $\det A$.}
    \item \question{Soit $u$ l'endomorphisme de $\R^3$ dont la matrice par
rapport \`a la base canonique est $A$. Pour quelles valeurs de $m$
est-ce un isomorphisme de $\R^3$\,?}
    \item \question{On pose $m = 1$. Trouver une base du noyau de $u$.}
\end{enumerate}
}
